% 第28章思考题28.1--28.14独立概念解析；候选，非正式章末练习。
\section*{第28章\quad 原子中的电子：思考题解析}
\paragraph{思考题 28.1}
对束缚电子，位置和动量满足
\(\Delta x\Delta p_x\ge\hbar/2\)。
若经典轨道同时给出任一时刻精确位置、
动量，就与量子态的概率描述不符。
氢原子定态由三维波函数或密度算符描述，
\(|\psi_{nlm}|^2\) 是概率密度，
\emph{不能}把其最大概率位置当成经典轨道。
Bohr轨道仅是在少数能谱上有效的历史模型。

\paragraph{思考题 28.2}
简并是不同线性独立量子态对应同一能量本征值。
非相对论纯库仑势氢原子
\(E_n=-13.6\,\mathrm{eV}/n^2\)，
对每个 \(l=0,\ldots,n-1\)，
\(m_l=-l,\ldots,l\)：
\[
\boxed{g_n=\sum_{l=0}^{n-1}(2l+1)=n^2
\quad\text{（不计自旋）}}.
\]
计入自旋且不考虑精细结构时是 \(2n^2\)；
外场、Lamb位移和精细结构会打破部分简并。

\paragraph{思考题 28.3}
多电子钾原子的价电子感受
核及其余电子的\emph{屏蔽}和\emph{交换}效应，
能量不仅依赖 \(n\)，也与 \(l\)、
电子关联和自旋轨道耦合有关。
相同主量子数 \(n\) 的不同 \(l\) 通常不再简并，
例如 \(4s\) 和 \(4p\)。
球对称场下的非相对论中心场近似，
\(m_l\) 不影响能量；
加入自旋轨道相互作用可出现总角动量 \(j\) 分裂。
题目“由哪些量子数决定”应注明使用哪个近似。

\paragraph{思考题 28.4}
在CPT对称以及相同外电磁场条件下，
氢与反氢的对应内部能级差相同，
因此
\[
\boxed{\nu_{\bar H}=\nu_H}
\]
（同一对应跃迁和物理校正下）。
反氢是反质子和正电子，
不应把两者电荷分别写成普通氢的。
若测得经充分消除环境效应的本征谱存在差别，
将提示基本对称性的检验问题；
不能简单将实验中不同外场磁矩取向造成的谱线
分裂认作CPT破坏。

\paragraph{思考题 28.5}
\(n=3\) 可有
\(l=0,1,2\)，分别记为
\(3s,3p,3d\)。按Pauli容量
\(2(2l+1)\)，
\[
\boxed{3s:2,\quad 3p:6,\quad 3d:10},
\qquad
\boxed{\text{合计 }18\text{ 个电子}}.
\]
这只是最多电子数，不代表按同一个 \(n\)
的全部轨道在多电子原子中能量相同。

\paragraph{思考题 28.6}
在半径 \(r\) 的经典库仑圆运动中，
向心条件
\(m_ev^2/r=e^2/(4\pi\epsilon_0r^2)\)，
故
\(K=m_ev^2/2=e^2/(8\pi\epsilon_0r)\)，
势能 \(V=-e^2/(4\pi\epsilon_0r)\)：
\[
\boxed{E_{\rm class}=K+V
=-\frac{e^2}{8\pi\epsilon_0r}}.
\]
若仅以连续经典理论分析，
任何稳定圆轨道半径都可形式上存在，
且带电加速电子应持续辐射。
Bohr的额外量子化假设
\(m_evr=n\hbar\) 才给
\(r_n=a_0n^2\) 和 \(E_n\propto-1/n^2\)；
这些不是纯经典力学的结论。

\paragraph{思考题 28.7}
经典磁偶极角可连续取向，
在非均匀磁场受力
\(F_z\simeq\mu_z\,\partial B/\partial z\)，
故经典预期是连续宽展的轨迹/斑带。
银原子基态的有效未成对电子
\(s=\frac12\)，对某一测量轴投影
\(m_s=\pm\frac12\)，使\(\mu_z\) 离散两值，
\[
\boxed{F_z\simeq\pm\mu_B\frac{\partial B}{\partial z}}.
\]
实验两束体现角动量投影量子化，
但真实束斑仍因入射热速度、狭缝和碰撞有宽度。

\paragraph{思考题 28.8}
氦原子基态 \(1s^2\)，两电子均
\[
\boxed{n=1,\ l=0,\ m_l=0},
\qquad
\boxed{m_{s,1}=+\tfrac12,\quad
m_{s,2}=-\tfrac12}.
\]
Pauli原理禁止两个全同电子\emph{全部四个}
量子数完全一致。
两电子自旋相反、成自旋单态，
空间部分在交换下对称，
全体费米波函数对交换反对称。

\paragraph{思考题 28.9}
连续轫致辐射的短波限由
\(hc/\lambda_{\min}=eU\) 决定，
\[
\boxed{\lambda_{\min}=\frac{hc}{eU}}
\]
与靶材无关；换靶且维持电压不变时
截止波长不变。
特征K\(\alpha\) 线与元素内层能级差有关。
铜 \(Z=29\) 小于银 \(Z=47\)，
依Moseley关系
\(E_{K\alpha}\propto(Z-\sigma)^2\)
会变小，因此
\(\boxed{\lambda_{K\alpha}(\mathrm{Cu})
>\lambda_{K\alpha}(\mathrm{Ag})}\)。
须先确认电压足以电离所讨论靶材的K壳层，
否则该特征谱线可能不出现。

\paragraph{思考题 28.10}
光子自旋为1，是\emph{玻色子}，
遵循Bose--Einstein统计，不服从对
\emph{相同单粒子状态的占据数}的Pauli排斥。
因此一个光学模式可以有任意个光子，
\[
\boxed{n_{\boldsymbol k,\lambda}=0,1,2,\ldots}.
\]
激光的受激辐射利用了该玻色增强；
但多光子数并不意味着可违反总能量或其他守恒律。

\paragraph{思考题 28.11}
两能级 \(E_2>E_1\) 的简化分布，
正常正温热平衡有 \(N_2/N_1
=(g_2/g_1)e^{-(E_2-E_1)/(k_BT)}\)。
若\emph{按每个态的平均占据}
\(N_2/g_2>N_1/g_1\)，
则可定义形式逆温
\[
\boxed{\beta_{\rm eff}
=-\frac{\ln[(N_2/g_2)/(N_1/g_1)]}{E_2-E_1}<0},
\quad T_{\rm eff}=(k_B\beta_{\rm eff})^{-1}<0.
\]
仅具有能量上限的有效子系统可出现这种负温度分布。
多能级激光器整体（包括平动自由度）不一定处于
负绝对温度的全局热平衡。

\paragraph{思考题 28.12}
空气—玻璃界面的Brewster角满足
\(\tan\theta_B=n_g/n_a\)。
在此角度，\emph{p偏振}
（电场在入射面内）菲涅耳反射系数为零，
而s偏振仍有非零反射。
两端斜置的Brewster窗因反复通过光学谐振腔，
使s偏振经历较大腔损耗，
最终优先振荡并输出
\[
\boxed{\boldsymbol E\text{ 位于入射面内的p线偏振}}.
\]
一个Brewster界面本身不意味着所有未知入射态
经过一次透射就严格变成纯p偏振；
谐振腔选择性损耗是关键。

\paragraph{思考题 28.13}
电子激发通常是eV量级，跃迁在可见、
近紫外乃至更短波段；
分子振动跃迁常约
\(10^{-2}\)--\(10^{-1}\,\mathrm{eV}\)，
多落红外；转动间距常约
\(10^{-4}\)--\(10^{-3}\,\mathrm{eV}\)，
多在远红外和微波。
分子的电子、振动、转动量子数同时改变，
产生紧密相邻的转动子线，
叠加在振动和电子谱带周围，
构成\emph{带状分子光谱}。
具体能标取决于分子的质量、键长和势能曲率。

\paragraph{思考题 28.14}
若相邻转动或振动能级差
\(\Delta E\gg k_BT\)，
从基态热跃迁的占据概率
\(\propto e^{-\Delta E/(k_BT)}\) 极小，
对比热的增量被指数压低，这就是
\emph{热冻结}。
当 \(k_BT\gtrsim\Delta E\)，
自由度逐渐恢复经典均分行为。
分子振动\emph{零点能}
\(E_0=\hbar\omega/2\) 虽仍存在，
但其对热容的温度导数为零。
电子能级常跨越几eV，
在室温 \(k_BT\simeq0.026\,\mathrm{eV}\)
通常冻结，电子激发对分子热容几乎无贡献；
金属的低能费米面电子激发是另一种情况。
